\subsection{Characterization of regular extensions}

PROPOSITION 9. --- Let K be a field, \(\bar{K}\) an algebraic closure of K and L an extension of K. Then the following conditions are equivalent:

\begin{enumerate}
\def\labelenumi{\alph{enumi})}
\item
  L is separable over K and K is relatively algebraically closed in L.
\item
  L is a regular extension of K.
\item
  The ring \(\bar{\mathbf{K}}\otimes_{\mathbb{K}}\mathbf{L}\) is an integral domain.
\item
  Let \(\bar{L}\) be an algebraic closure of L. Then L is linearly disjoint over K from the relative algebraic closure of K in \(\bar{L}\) .
\end{enumerate}

Moreover, when these conditions hold, then \(\bar{\mathbf{K}}\otimes_{\mathbf{K}}\mathbf{L}\) is a field.

\(a) \Rightarrow b)\) : Let M be an extension of K. Under the hypotheses of \(a\) ), the ring M \(\otimes_{K} L\) is an integral domain, by V, p.~140, Cor.

\(b) \Rightarrow c)\) : This follows from Def. 2.

\(c) \Rightarrow d): \text{With the notation of } d) \text{ we can identify } \bar{\mathbf{K}} \text{ with the relative algebraic closure of } \mathbf{K} \text{ in } \bar{\mathbf{L}} (\mathbf{V}, \text{ p. 22, Ex. 2}). \text{ Suppose that the ring } \mathbf{A} = \bar{\mathbf{K}} \otimes_{\mathbf{K}} \mathbf{L} \text{ is an integral domain. Let } \mathbf{E} \text{ be a subextension of } \bar{\mathbf{K}} \text{ of finite degree over } \mathbf{K}; \text{ the subring}\)

\(E \otimes_{K} L\) of A is an integral domain and hence is an algebra of finite degree over L; by the Cor. of V, p.~10 it is a field. Since \(\bar{K}\) is the union the increasing directed set of extensions E of the above type, A is a field (V, p.~11, Prop. 3). The canonical homomorphism of A into \(\bar{L}\) mapping \(x \otimes y\) to xy (for \(x \in \bar{K}\) and \(y \in L\) ) is therefore injective, so L and \(\bar{K}\) are linearly disjoint over K.

\(d) \Rightarrow a): \text{Under the hypotheses of } d) \text{ we have } L \cap \bar{K} = K, \text{ so } K \text{ is relatively algebraically closed in } L; \text{ further if } p \text{ is the characteristic exponent of } K, \text{ the field } L \text{ is linearly disjoint from } K^{p^{-\infty}} \text{ over } K, \text{ hence } L \text{ is separable over } K (V, p. 123, Cor. 1).\)

COROLLARY 1. --- Let A be an algebra over a field K. For A to be a regular K-algebra it is necessary and sufficient for the ring \(\bar{\mathbf{K}}\otimes_{\mathbf{K}}\mathbf{A}\) to be an integral domain.

The stated condition is clearly necessary. Conversely, assume that \(\bar{K} \otimes_{K} A\) is an integral domain and denote by E the field of fractions of A. By Prop. 4 (V, p.~141) the ring \(\bar{K} \otimes_{K} E\) is an integral domain, hence E is a regular extension of K, by Prop. 9; by V, p.~141, Cor., we conclude that A is regular as K-algebra.

COROLLARY 2. --- Let K be an algebraically closed field. Every K-algebra which is an integral domain is a regular K-algebra. In particular, every extension of K is regular.

This follows from Cor. 1.

COROLLARY 3. --- Let K be an algebraically closed field. If A and B are two K-algebras which are integral domains, then the same is true of A ⊗\_K B.

By Cor. 2, A and B are regular K-algebras, and it suffices to apply Prop. 2 (V, p.~141).

